% ============================================================
% Plantilla: Artículo IEEE para revistas (clase IEEEtran, modo journal)
% Citas numéricas: \cite{clave} -> [1]. Compila con pdflatex + bibtex.
% ============================================================
\documentclass[journal]{IEEEtran}

\usepackage[utf8]{inputenc}
\usepackage[T1]{fontenc}
\usepackage[spanish,es-tabla]{babel}
\usepackage{amsmath,amssymb}
\usepackage{graphicx}
\usepackage{booktabs}

% IEEEtran no traduce este título con babel
\renewcommand{\IEEEkeywordsname}{Palabras clave}

\begin{document}

\title{Título del Artículo para la Revista IEEE}

\author{Primer~Autor,~\IEEEmembership{Miembro,~IEEE,}
  y~Segundo~Autor%
  \thanks{P. Autor pertenece al Departamento de Ingeniería, Universidad, Ciudad, País
    (correo: autor@ejemplo.com).}%
  \thanks{Manuscrito recibido el 1 de enero de 2026.}}

\markboth{Revista IEEE de Ejemplo,~Vol.~1, No.~1, enero~2026}%
{Autor \MakeLowercase{\textit{et al.}}: Título abreviado del artículo}

\maketitle

\begin{abstract}
Resume en un solo párrafo el problema, el método propuesto, los resultados principales
y su relevancia (máximo 250 palabras).
\end{abstract}

\begin{IEEEkeywords}
Palabra clave uno, palabra clave dos, palabra clave tres.
\end{IEEEkeywords}

\section{Introducción}
\IEEEPARstart{E}{sta} plantilla sigue el formato de las revistas IEEE. La primera letra
del artículo se destaca con \texttt{\textbackslash IEEEPARstart}. Las referencias se
numeran en orden de aparición \cite{lamport1994}.

\section{Trabajos relacionados}
Resume el estado del arte \cite{knuth1984,goossens2004}.

\section{Método propuesto}
El modelo se define como:
\begin{equation}
  y = \mathbf{W}x + b
  \label{eq:modelo}
\end{equation}
donde $\mathbf{W}$ es la matriz de pesos y $b$ el sesgo.

\section{Resultados}
La Tabla~\ref{tab:resultados} compara el método propuesto, basado en~\eqref{eq:modelo}.

\begin{table}[!t]
  \caption{Comparación de métodos}
  \label{tab:resultados}
  \centering
  \begin{tabular}{lcc}
    \toprule
    Método & Exactitud & F1 \\
    \midrule
    Referencia & 0.84 & 0.81 \\
    Propuesto & \textbf{0.91} & \textbf{0.89} \\
    \bottomrule
  \end{tabular}
\end{table}

\section{Conclusiones}
Resume las contribuciones y el trabajo futuro.

\bibliographystyle{IEEEtran}
\bibliography{referencias}

\begin{IEEEbiographynophoto}{Primer Autor}
recibió el título de Ingeniero en Software de la Universidad en 2026. Sus intereses de
investigación incluyen la seguridad de la información.
\end{IEEEbiographynophoto}

\end{document}
